% =====================================================================
%  Chapter 13 — Unification (S36–S37)
% =====================================================================
\chapter{Unification: Two Open Problems Become One}
\label{ch:unification}

\begin{record}
\textbf{What this chapter covers.} Sessions \Sn{36} and \Sn{37}, run
back-to-back on the same day. \Sn{36} verifies the two couplings the previous
session had left pending, retracts a term that session had claimed, and
isolates the one residue that survives. \Sn{37} derives H2 from H1 and
discovers that in the generic case H2's remaining content and the
\(\sstar\)-decay gap are \emph{the same open problem}.

\smallskip
That discovery is the programme's cleanest structural result. It reduced three
nominally independent open items to two, and it did so without proving
anything new about either. It is worth being precise about what such a result
is: a reduction in the number of things that must be proved, obtained by
showing that two of them were always one. Nothing became easier. The map
became accurate.
\end{record}

\section{\texorpdfstring{\Sn{36}: verification, retraction, residue}{S36: verification, retraction, residue}}
\label{sec:s36}

\Sn{35} had recorded \texttt{prop:K-self} (line~6026), the
\(K\)-self-absorption identity, explicitly marked \emph{provisional}, with two
couplings flagged for verification:

\begin{itemize}[leftmargin=1.6em]
\item \textbf{C1}, the \(\Dt w\) substitution, including a
  \(\nabla\log m\cdot\nabla w\) re-entry whose constant had to come out below
  one;
\item \textbf{C2}, the \(\alpha\)-weighted marginal terms, a constant chase
  with logarithmic degradation.
\end{itemize}

If both verified, the consequence stated in advance was substantial:
\texttt{conj:K-refined} --- the \(c_K\) hypothesis --- would become
\emph{unnecessary}, and H1 would reduce to \(\sstar\le\varepsilon_0\) plus
routine items.

\subsection{C2, and then C1}

\textbf{C2 verified.} The \(\alpha\)-weighted terms close via the a priori
bound of \texttt{thm:apriori-sigma} (line~5873), at the cost of an
\(\Lfac^2\) degradation --- which lands in the class the estimate already
allows.

\textbf{C1 verified, modulo one item, and the mechanism is better than
required.} The dangerous couplings are the drift--curvature terms, and they do
not merely turn out to be bounded: they are \emph{good-signed}.

\begin{statement}[Gram negativity --- \texttt{lem:gram-negativity}, line~6180]
\label{st:gram}
\PROVED{} The drift--curvature terms contract \(\nabla m\otimes\nabla m\)
against the Gram tensor \(G_{jk}=\partial_j\ehat\cdot\partial_k\ehat\) with a
\emph{negative} sign. Since \(G\) is positive semi-definite by construction,
the most dangerous couplings help rather than merely fail to hurt.
\end{statement}

This is the second place in the programme where
\texttt{lem:grad-omega-split} --- the audit's survivor --- is what makes the
argument work: it is the identity that makes \(G\) the natural object and the
negativity a statement about a positive semi-definite tensor rather than a
computation with signs to be tracked term by term.

The stretching and commutator couplings close via the Biot--Savart near/far
split already in place. One item resisted: the transition-annulus terms,
supported on \(\{M/8\le\abs{\omega}\le M/4\}\), which appeared to require the
standard De Giorgi nested-level iteration. \Sn{36} classified that as
\emph{open-mechanical} --- ``classical bookkeeping, distinct in kind from the
genuinely open items''.

That classification was wrong, and Chapter~\ref{ch:annulus} is the session
that proved it wrong.

\subsection{The retraction}

\Sn{35}'s identity had claimed \emph{two} good terms on the left: the
\(K\)-density \(\nu\iint\abs{\nabla m}^2w\eta^2\), and additionally
\(\nu\iint m^2w^2\eta^2\) --- the critical quadratic with a favourable weight,
presented as reinforcing the second rung.

\begin{statement}[\texttt{lem:w2-cancel}, line~6166]
\label{st:w2cancel}
\WITHDRAWN{} The \(+\nu w\ehat\) term of the direction equation contributes
exactly \(+\nu w^2\) to \(\tfrac12\Dt w\), which cancels the sink
\(-\nu\iint m^2w^2\eta^2\) \emph{identically}. The second left-hand term does
not exist.
\end{statement}

The \(K\)-control conclusion survives with amended constants, and no
downstream conclusion had depended on the cancelled term --- the critical
quadratic is handled solely by \texttt{prop:second-rung}, which suffices. The
episode is nonetheless recorded as Wall~7, and the walls document is explicit
about why: it is one of three instances of the same failure mode, ``a
favourable statement written down before the computation producing it was
closed''. The other two are the \Sn{34} corrigendum and the inverted lemma of
\Sn{44} (Chapter~\ref{ch:non-implication}).

Two things are worth noting about the retraction. It was found by the
programme auditing its own previous session before building on it, which is
the rule of \S\ref{sec:reporting-rules} in Chapter~\ref{ch:house-rules}. And
the provisional label was on the proposition from the start: \Sn{35} wrote
down what it had, marked what it had not checked, and named the two items ---
which is why \Sn{36} knew where to look.

\subsection{What \texorpdfstring{\Sn{36}}{S36} left standing}

The routine items (i) transport, (ii) commutator and (iv) time slice were
discharged with short proofs. \texttt{prop:K-amended} (line~6215) then bounds
the \(K\)-coupling with an \emph{absolute constant}, no logarithm and no
correlation hypothesis --- modulo one thing.

\begin{defn}[The annulus residual --- \texttt{def:annulus-residual}, line~6234]
\label{def:aann}
\(\Aann\) consists of the cutoff-derivative terms supported on the transition
annulus \(\mathrm{Ann}=\{M/8\le\abs{\omega}\le M/4\}\cap\operatorname{supp}\psi\):
\[
  \nu\iint_{\mathrm{Ann}}\abs{\nabla m}^2\,w\,\psi\abs{\chi'},
  \qquad
  \nu\iint_{\mathrm{Ann}}m^2w^2\,\psi\abs{\chi'},
\]
and lower-order variants, all carrying \emph{incomplete} cutoff weights
\(\psi\abs{\chi'}\) rather than \(\eta^2\).
\end{defn}

\texttt{rem:level-iteration} (line~6249) sketched the disposal: replace the
single level pair by a nested sequence \(\theta_k=\tfrac18+2^{-k-3}\) with
Young parameters \(\delta_k=\delta_08^{-k}\), on the principle that
``geometric beats geometric''. The remark ends: ``This is bookkeeping of the
same class as Ladyzhenskaya--Ural'tseva Chapter~II; it involves no new
estimate.''

The remark now carries a red superseded block, added at \Sn{41}. The scheme
was executed and does not close.

\begin{table}[ht]
\centering
\small
\begin{tabular}{@{}p{0.34\textwidth} p{0.28\textwidth} p{0.28\textwidth}@{}}
\toprule
Component of H1 & Status after \Sn{36} & Status today \\
\midrule
First rung (angular Caccioppoli) & Proved; items (i),(ii),(iv) discharged
  & \PROVED \\[2pt]
Critical quadratic & Proved given \(\sstar\le\varepsilon_0\)
  & \CONDITIONAL{\(\varepsilon_0\)} \\[2pt]
\(K\)-absorption & Proved modulo \(\Aann\)
  & \PROVEDMOD{\(\Aann\)} \\[2pt]
C2 (\(\alpha\)-terms) & Verified & \PROVED \\[2pt]
Gram negativity & Proved, good signs & \PROVED \\[2pt]
\(w^2\) cancellation & Corrects \Sn{35} & \WITHDRAWN{} (the claimed term) \\[2pt]
\(\Aann\): families G1--G4 & --- & \PROVED{} good-signed, discarded (\Sn{41}) \\[2pt]
\(\Aann\): families I1--I2 & --- & \PROVED{} in the allowed class (\Sn{41}) \\[2pt]
\(\Aann\): core S1\,+\,S2 & ``open-mechanical''
  & \OPEN{}; level iteration \IMPOSSIBLE{} \\
\bottomrule
\end{tabular}
\caption[H1 after \Sn{36}, and now]{The consolidated status of H1's estimate
layer, as \Sn{36} recorded it and as it stands. The only row that changed is
the last, and it changed for the worse: what \Sn{36} believed was free
bookkeeping is genuine analytical content.}
\label{tab:h1-status}
\end{table}

\subsection{The depletion diagnostic}

\Sn{36} also built the first numerical instrument aimed at H2:
\texttt{layer4/depletion\_diagnostic.py}, measuring \(\alpha(x^*)/M\), the
quantity H2 needs to be small and falling.

\begin{center}
\small
\begin{tabular}{@{}l r r r r@{}}
\toprule
\texttt{exp\_id} & median ratio & max & final & \(M\) range \\
\midrule
\texttt{EXP-L4-DEPLETION-TG-001} & \(0.231\) & \(0.538\) & \(0.271\)
  & \([1.40,\,20.94]\) \\
\texttt{EXP-L4-DEPLETION-ADV-001} & \(0.016\) & \(0.032\) & \(0.021\)
  & \([1.75,\,1.96]\) \\
\bottomrule
\end{tabular}
\end{center}

\medskip
The Taylor--Green timeseries anti-correlates the ratio with \(M\): it is
\(0.538\) at \(M=1.40\), near the minimum of \(M\), and falls to \(0.085\) at
\(M=18.4\), near the peak. The recorded \(\theta_{\mathrm{eff}}\) --- the
effective scaling exponent of \(\alpha\) in \(M\), defined only for \(M>10\)
--- has median \(-0.080\), which puts the observed exponent near \(0.42\),
well below the naive value \(1\). Read as a statement about ordinary vortex
dynamics, this is evidence of genuine depletion. Read as evidence about H2 it
is worth much less, for the reason Wall~6 gives and
Chapter~\ref{ch:numerical-lessons} presses: the flow is globally regular and
decaying.

\section{\texorpdfstring{\Sn{37}: deriving H2, and finding it is \(\sstar\)-decay}{S37: deriving H2, and finding it is sigma-star-decay}}
\label{sec:s37}

H2 was stated at \Sn{32} as an open target with a named literature to adapt:
Constantin--Fefferman assumes coherence at a \emph{fixed} scale \(\rho\),
whereas the pincer needs it at the \emph{shrinking} scale \(\rstar(t)\).
\Sn{37} carried out the adaptation.

\subsection{The derivation}

\texttt{prop:localized-cf} (line~6381) is classical technique applied to the
programme's geometry. Constantin's representation of the stretching rate
carries a geometric kernel with an exact pointwise cancellation,
\(D(e,e,\sigma)=0\). Lipschitz continuity of \(\ehat\) --- which is exactly
what H1's Bridge Lemma delivers --- controls the near field; Cauchy--Schwarz
against the global enstrophy \(\Omega=\norm{\omega}_{L^2}^2\) controls the far
field. Optimizing the cutoff radius between the two gives
\[
  \abs{\alpha(x^*)} \;\lesssim\; (KM)^{3/5}\,\Omega^{1/5},
  \qquad K = \frac{C_B(\sstar)^{1/2}}{\rstar} = C_B(\sstar)^{1/2}M^{1/2},
\]
and substituting \(K\),
\begin{equation}\label{eq:cf-master}
  \abs{\alpha(x^*)} \;\le\; C\,(\sstar)^{3/10}\,M^{9/10}\,\Omega^{1/5}.
\end{equation}

\begin{record}
\textbf{A self-audit inside the derivation.} The session's first-pass sketch
assumed \(\Omega=O(1)\) for free. That is false near a putative singularity:
for an exactly self-similar Type-I profile,
\(\omega(x,t)\sim(T-t)^{-1}\Omega_0\!\left((x-x^*)/\sqrt{T-t}\right)\) gives
\(\norm{\omega(\cdot,t)}_{L^2}^2\sim(T-t)^{-1/2}\sim M(t)^{1/2}\), not
\(O(1)\). Assuming bounded enstrophy near a singularity is close to assuming
there is no singularity. The derivation was redone with the correct scaling,
and that redoing is what produced the dichotomy below.
\end{record}

\subsection{The dichotomy}

Everything now turns on how \(\Omega\) grows.

\begin{statement}[Regime A --- \texttt{prop:regime-A}, line~6430]
\label{st:regA}
If \(\Omega(t)=O(1)\) then \eqref{eq:cf-master} gives
\(\abs{\alpha(x^*)}\le C(\sstar)^{3/10}M^{9/10}\), i.e.\ \(\theta=2/5<1/2\):
\textbf{H2 holds unconditionally given H1}, using only the a priori bound
\(\sstar\le C\Lfac/\nu\), with no decay required.
\end{statement}

\begin{statement}[Regime B --- \texttt{prop:regime-B}, line~6442]
\label{st:regB}
If \(\Omega(t)\sim M(t)^{1/2}\) --- the scaling of a genuine self-similar
Type-I profile --- then
\(\abs{\alpha(x^*)}\le C(\sstar)^{3/10}M^{9/10}\cdot M^{1/10}
= C(\sstar)^{3/10}M\), which is \textbf{exactly borderline},
\(\theta=1/2\). The near/far optimization gains nothing beyond the constant
prefactor \((\sstar)^{3/10}\). Strict depletion requires \(\sstar(t)\to0\),
not merely \(\sstar\) bounded.
\end{statement}

\begin{figure}[ht]
\centering
\begin{tikzpicture}[
  >={Latex[length=2mm]},
  bx/.style={draw, rounded corners=2pt, font=\scriptsize, align=center,
             inner sep=3.5pt, text width=37mm, minimum height=10mm},
  gd/.style={bx, draw=provedgreen, thick},
  bd/.style={bx, draw=impossiblered, thick},
  nx/.style={bx, draw=openblue, thick, fill=panelgrey},
  ar/.style={->, thick}
]
\node[bx] (cf) at (0,0)
  {\eqref{eq:cf-master}:\\
   \(\abs{\alpha}\lesssim(\sstar)^{3/10}M^{9/10}\Omega^{1/5}\)};

\node[gd] (A) at (5.6,1.6)
  {\textbf{Regime A}: \(\Omega=O(1)\)\\
   \(\theta=2/5<\tfrac12\)\\ H2 unconditional given H1};
\node[nx] (B) at (5.6,-1.6)
  {\textbf{Regime B}: \(\Omega\sim M^{1/2}\)\\
   \(\theta=\tfrac12\) exactly\\ needs \(\sstar\to0\)};

\node[bd] (Aimp) at (11.0,1.6)
  {\Sn{38}: \IMPOSSIBLE\\
   \(\limsup\Omega=\infty\) at any\\ genuine blow-up};
\node[nx] (unif) at (11.0,-1.6)
  {\(\Rightarrow\) H2 \(\equiv\) \(\sstar\)-decay\\
   \emph{one} open problem, not two};

\draw[ar] (cf) -- (A);
\draw[ar] (cf) -- (B);
\draw[ar, impossiblered] (A) -- (Aimp);
\draw[ar] (B) -- (unif);

\node[font=\scriptsize\itshape, text=rulegrey, align=center, text width=95mm]
  at (5.6,-3.4)
  {\Sn{37} produced the dichotomy and unified the lower branch. \Sn{38}
   removed the upper branch. What remains is a single statement, and it is the
   statement \Sn{30} once believed it had proved.};
\end{tikzpicture}
\caption[The Regime A/B dichotomy and its collapse]{The two regimes of
\texttt{prop:regime-A} and \texttt{prop:regime-B} (lines 6430, 6442), and what
happened to each. The programme entered \Sn{37} with three open items and left
\Sn{38} with one deep one --- by proving that one branch was already the
other, and that the remaining branch does not exist.}
\label{fig:regimes}
\end{figure}

\subsection{The unification}

\begin{statement}[\texttt{cor:unification}, line~6456]
\label{st:unification}
In the generic case --- Regime~B or worse --- \textbf{H2's remaining content
and the \(\sstar\)-decay gap are the same open problem}. Both require
\(\sstar(t)\to0\), not merely the a priori bound. The programme's genuinely
open analytical content therefore reduces from three nominally independent
items (H1's annulus step, H2, \(\sstar\)-decay) to two: the annulus
level-iteration, then thought mechanical, and a single deep gap.
\end{statement}

The corollary as written keeps Regime~A available: it notes that a genuine
subcritical enstrophy bound near the singularity ``remains an interesting
alternative target, since it does not require touching \(\sstar\) at all''.
That sentence was true for one session. \Sn{38} proved Regime~A impossible
(Chapter~\ref{ch:pincer}, \S\ref{sec:regime-a}), and with it the only
alternative.

\subsection{The enstrophy exponent, measured}

The natural diagnostic follows immediately: which regime does a real flow sit
in? \Sn{37} built \texttt{layer4/enstrophy\_exponent.py} and fitted \(p\) in
\(\Omega\sim M^p\) on the Taylor--Green stretching run, chosen with the same
seed and parameters as \texttt{EXP-L4-DEPLETION-TG-001} for direct
comparability.

\begin{center}
\small
\begin{tabular}{@{}l l r r r@{}}
\toprule
\texttt{exp\_id} & fit window & \(p\) & \(r^2\) & points \\
\midrule
\texttt{EXP-L4-ENSTROPHY-EXPONENT-TG-001} & all, \(M>5\) & \(1.226\) & \(0.624\) & \(33\) \\
\texttt{EXP-L4-ENSTROPHY-EXPONENT-TG-001} & growth phase only & \(\mathbf{0.891}\) & \(\mathbf{0.914}\) & \(19\) \\
\texttt{EXP-L4-ENSTROPHY-EXPONENT-ADV-001} & fallback, \(M>3\) & \(16.088\) & \(1.000\) & \(\mathbf{2}\) \\
\bottomrule
\end{tabular}
\end{center}

\medskip
The Taylor--Green growth-phase fit, \(p\approx0.89\) with \(r^2=0.91\), lands
\emph{above} both Regime~A (\(p\approx0\)) and Regime~B (\(p=1/2\)) ---
nominally worse than borderline. The session's own reading is the correct one
and is not softened: this measures ordinary transient vortex stretching in a
globally regular, \emph{decaying} flow whose \(M\) peaks and then falls, not
the approach to a singularity. It is mild evidence against expecting bounded
enstrophy as a free lunch even in ordinary dynamics, which sharpens the case
that \(\sstar\)-decay is the unavoidable core difficulty. It does not test the
theory.

The adversarial row is the more instructive one for this book's purposes. Its
\(M\) range is \([1.746,\,1.962]\) --- essentially flat --- so the fallback
threshold admitted exactly two points, and a two-point fit returns
\(r^2=1.000\) by arithmetic. The session log says the run was ``correctly
reported as NaN/2-point artifact rather than forced''. The database row
disagrees: it records \(p=16.088\), \(r^2=1.0000\), \(n=2\), and
\texttt{verdict = PASS}, with a \texttt{note} field recording only the
threshold fallback.\footnote{Source disagreement, resolved in favour of the
database, which is the higher authority for experimental claims
(Chapter~\ref{ch:how-to-read}, \S\ref{sec:source-hierarchy}). The narrative in
\texttt{SESSION-LOG/2026-07-19-S37-cf-depletion-unification.md} is what the
session intended; the row is what was written. A \textsf{PASS} on a two-point
regression is not a result, and the row should be read as vacuous. It is
listed in Appendix~\ref{ch:experiments} with that reading attached.}

\begin{obsv}[What a unification is worth]
\label{obs:unification-value}
\Sn{37} is the session the programme is most likely to be proud of, and it is
worth being exact about what it achieved. It proved no new estimate about
\(\sstar\). It made no target easier. What it did was show that two of the
programme's three open items had always been one, so that effort spent on
either was effort spent on both --- and, in the same stroke, that the more
attractive of the two branches (Regime A, unconditional depletion, no decay
required) was the one that depended on an assumption nobody had checked.
\Sn{38} then checked it.

\smallskip
The pattern is worth naming because it recurs: in this programme, the sessions
that improved the \emph{map} were more valuable than the sessions that
extended the \emph{route}, and they were also the sessions that most often
delivered bad news.
\end{obsv}
